%AttackName: M-MGA-1

%Input:
The input is the matrix \(A\) and the objective function \(\hat{L}_t(\hat{A})\) to be maximized.

%Output:
The output is the optimized matrix \(\hat{A}\), which maximizes the objective function while satisfying the constraints.

%Formula:
$\arg \max_{\hat{A}} \hat{L}_t(\hat{A}), \quad \text{s.t.} \quad i < j, \quad |A_{ij} - \hat{A}_{ij}| \leq \rho$

%Explanation:
s:
Matrix Based Method MGA where the objective function is given by:
$\hat{L}_t(\hat{A})$
This function represents the objective to be maximized. It depends on the matrix \(\hat{A}\) and is typically designed to encourage misclassification or decrease confidence in the correct classification of the input data.
optimization variable:
The optimization variable is:
$\hat{A}$
This variable represents the matrix that is being optimized. The goal is to find the values of $\hat{A}$ that maximize the objective function $(\hat{L}_t(\hat{A})$
constraints:
The constraints are defined as:
$i < j$
and
$|A_{ij} - \hat{A}_{ij}| \leq \rho$
The first constraint ensures that the elements $i$ and $j$ of the matrix $\hat{A}$ are distinct, where $i$vand $j$ are indices representing different features or components. The second constraint limits the allowable change between the original matrix $A$ and the optimized matrix $\hat{A}$ to be less than or equal to $\rho$. This constraint controls the magnitude of the perturbations applied to the matrix elements.